\section{Method}

\subsection{Overall Pipeline}

The proposed pipeline processes the input video frame by frame.
First, representative frames are selected from the video.
Then, the piano keyboard is localized and isolated.
White and black keys are detected on the keyboard region.
Fingertips are detected on the full image.
Finally, fingertip positions are associated with keys
to identify pressed keys at each time instant.

\subsection{Frame Selection}

From each input video, a subset of frames is extracted.
Frames with minimal hand occlusion are selected.
These frames are used for keyboard and key detection.

\subsection{Keyboard Localization}

The piano keyboard is localized in the selected frames.
The keyboard region is isolated from the background
to simplify further processing.

\subsection{Piano Key Detection}

White and black piano keys are detected inside the keyboard region.
White keys are detected first to define the keyboard structure.
Black keys are detected relative to the white key layout.

\subsection{Fingertip Detection}

Fingertips are detected on the full image.
Detection is performed independently from key detection.
This ensures that fingertip positions remain consistent
with the original image coordinates.

\subsection{Key-Finger Association}

Detected fingertip positions are compared with key locations.
A fingertip is assigned to a key if it lies inside the key region.
Temporal consistency is used to reduce false detections.

\subsection{Pressed-Key Detection}

The pressed keys are determined through the combination of spatial and motion-based information. First, the spatial information of the position of the fingertips is mapped to the identified key regions. A key is said to be pressed if the fingertip is spatially associated with the key and the finger corresponding to the key is identified as pressed through articulation analysis.
